% Optional replacement for the analytical example in the theory section.
% Keep this figure outside the Simulation section to avoid duplication.
\providecommand{\simroot}{.}
To compare the complexity structures at the same total coverage, fix
$D=2$, $\sigma^2=\sigma_{\rm tr}^2=1/4$, and $\delta=0.05$.
For the separate bound, equal budgets are
$\epsilon_{\rm ch}=\epsilon_{\rm data}=1-\sqrt{1-\delta}$.
The adaptive allocation has
$\epsilon_{\rm data}/\epsilon_{\rm ch}=99$, with both allocations
satisfying $(1-\epsilon_{\rm data})(1-\epsilon_{\rm ch})=1-\delta$.
Figure~\ref{fig:bound_comparison} evaluates the resulting expressions
with fixed $n=10$ and with $n=m$. In the asymmetric case, the joint
complexity crosses the separate complexity at approximately $m=396$
under equal allocation and $m=348$ under adaptive allocation. The
separate bounds approach data-complexity floors of $0.6159$ and
$0.5848$, respectively, whereas the joint logarithm keeps growing.
When $n=m$, the joint formulation is tighter over the illustrated range.
This comparison isolates the analytical bound structure and is not a
learning-curve experiment.
\begin{figure}[t]
\centering
\includegraphics[width=\linewidth]{\simroot/assets/fig_bound_structure.pdf}
\caption{Analytical complexity comparison with exact product-confidence
budget allocation, $D=2$ and $\sigma^2=\sigma_{\rm tr}^2=1/4$.}
\label{fig:bound_comparison}
\end{figure}
